\section{Appendix D: Testing bounds and minimax proofs}

The distance comparisons in \cref{obj:lem:baseline-translation-affinity,obj:lem:walsh-channel-energy,obj:lem:hidden-allocation-contraction} support a uniform testing comparison for the complete-record block experiment. The constant-target property of \cref{obj:lem:block-law} then connects testing difficulty to squared estimation risk. Together with the supplied-graph comparison and the observable upper guarantee, these results characterize attainable precision throughout the admissible degree and retention ranges.

\subsection{Uniform block testing and its radius}

For probability laws \(P\) and \(Q\) on a common measurable space, total-variation distance is
\[
\operatorname{TV}(P,Q)
=\sup_{A\text{ measurable}}|P(A)-Q(A)|.
\]
This convention agrees with \cref{obj:lem:baseline-translation-affinity}. The comparison concerns the opposite-sign complete-record mixture laws \(P_{+1,h}\) and \(P_{-1,h}\) of \cref{obj:def:block-prior}, where \(h\) is the block amplitude.

With \(m=Bd\), the source count, and \(p=1-(1-q)^d\), the association-revelation probability specified in \cref{obj:def:testing-handle}, define the testing scale by
\[
s(B,d,q)=
\begin{cases}
\min\{1,d/\sqrt{mp}\},&p>0,\\
1,&p=0.
\end{cases}
\]
The testing radius is the supremum of amplitudes in \([0,1/4]\) for which the two mixture laws have total variation at most one half:
\[
\bar h(B,d,q)
=\sup\left\{h\in[0,1/4]:
\operatorname{TV}(P_{+1,h},P_{-1,h})\le\tfrac12\right\}.
\]
The universal amplitude constant is \(\kappa=(100\pi)^{-1}\).

The preceding distance bounds have complementary roles in this comparison. Baseline translations control changes in response locations, while the Walsh energy bound measures response dependence on source signs. The hidden-allocation comparison accounts for the source associations remaining uncertain after the audit. Through the reconstruction property of \cref{obj:lem:block-law}, these comparisons concern the complete graph, assignment, and outcome record, including the observed global hidden treated-source count. The resulting statement provides both a uniformly small-distance range and a complementary lower bound on distinguishability.

\begin{theoremv}{thm:block-testing-scale}[Block testing scale bounds]
Suppose the block experiment satisfies the following conditions:
\begin{itemize}
\item \textbf{(Population and blocks.)} The population size \(n\), block count \(B\), and degree \(d\) are integers satisfying
\[
n\ge4,\qquad B\ge1,\qquad d\ge1,\qquad 2Bd\le n.
\]
\item \textbf{(Retention.)} The edge-retention probability satisfies \(q\in[0,1]\).
\item \textbf{(Assignment.)} The experimental design satisfies \cref{obj:ass:assignment}.
\item \textbf{(Audit.)} The experimental design satisfies \cref{obj:ass:audit} at retention probability \(q\).
\item \textbf{(Independence.)} The experimental design satisfies \cref{obj:ass:design-independence}.
\end{itemize}
Put \(m=Bd\), the number of sources, and \(p=1-(1-q)^d\), the probability that a source association is revealed. Define the testing scale and universal constant by
\[
s(B,d,q)=
\begin{cases}
\min\{1,d/\sqrt{mp}\},&p>0,\\
1,&p=0,
\end{cases}
\qquad
\kappa=\frac{1}{100\pi}.
\]
Let \(P_{+1,h}\) and \(P_{-1,h}\) be the complete original-record mixture laws under the opposite-sign block priors, using this experimental design, and let \(\operatorname{TV}\) denote total-variation distance. For every real amplitude \(h\),
\[
0\le h\le\kappa s(B,d,q)
\quad\Longrightarrow\quad
\operatorname{TV}(P_{+1,h},P_{-1,h})\le\frac12.
\]
For every real amplitude \(h\) satisfying \(0<h\le1/4\) and \(p>0\),
\[
\operatorname{TV}(P_{+1,h},P_{-1,h})
\ge 1-\frac{7d^2}{8h^2mp}.
\]
Finally, define the testing radius by
\[
\bar h(B,d,q)
=
\sup\left\{
h\in[0,1/4]:
\operatorname{TV}(P_{+1,h},P_{-1,h})\le\frac12
\right\}.
\]
Then
\[
\kappa s(B,d,q)\le\bar h(B,d,q)\le2s(B,d,q).
\]
\end{theoremv}

The scale in \cref{obj:thm:block-testing-scale} balances the amplitude against the number of revealed source associations. At positive revelation probability, the distance lower bound depends on the squared amplitude through \(h^2mp/d^2\), while the uniform upper bound supplies a range of amplitudes at which the signs remain difficult to distinguish. The truncation at one keeps the scale bounded when revelation is sparse. At zero revelation, the constant-amplitude comparison continues to apply. The radius bounds characterize the supremum directly, using comparisons that hold throughout their stated amplitude ranges.

Under the canonical assignment--audit product design, the same comparison can be recorded entirely in terms of the block parameters. The following statement consolidates the testing radius and the uniform amplitude guarantee, and gives the explicit degree-two expression.

\begin{theoremv}{thm:nonlocal-testing-answer}[Nonlocal block testing bounds]
Let the population size \(n\), block count \(B\), degree \(d\), and retention probability \(q\) satisfy:
\begin{itemize}
\item \textbf{(Population.)} \(n\) is an integer with \(n\ge4\).
\item \textbf{(Blocks and degree.)} \(B,d\) are integers with \(B\ge1\) and \(d\ge1\).
\item \textbf{(Capacity.)} \(2Bd\le n\).
\item \textbf{(Retention.)} \(q\in[0,1]\).
\end{itemize}
Define the source count \(m\), association-revelation probability \(p\), and testing scale \(s\) by
\[
m=Bd,\qquad p=1-(1-q)^d,\qquad
s=
\begin{cases}
\min\{1,d/\sqrt{mp}\},&p>0,\\
1,&p=0.
\end{cases}
\]
Let \(P_{+1,h}\) and \(P_{-1,h}\) be the complete original-record block mixture laws under the canonical assignment-audit product design with retention probability \(q\), at amplitude \(h\) and signs \(+1\) and \(-1\), respectively. With \(\operatorname{TV}\) denoting total-variation distance, define the testing radius
\[
\bar h
=\sup\left\{h\in[0,1/4]:
\operatorname{TV}(P_{+1,h},P_{-1,h})\le\tfrac12\right\}.
\]
Then
\[
(100\pi)^{-1}s\le\bar h\le2s.
\]
Moreover, for every real \(h\) satisfying \(0\le h\le(100\pi)^{-1}s\),
\[
\operatorname{TV}(P_{+1,h},P_{-1,h})\le\tfrac12.
\]
If \(d=2\), the scale satisfies
\[
s=
\begin{cases}
\min\left\{1,\sqrt{\dfrac{2}{B(2q-q^2)}}\right\},&q>0,\\
1,&q=0.
\end{cases}
\]
\end{theoremv}

For the canonical design, \cref{obj:thm:nonlocal-testing-answer} expresses testing difficulty through block count, degree, and retention. In the degree-two case, \(2q-q^2\) is the probability that at least one of a source's two outgoing arrows is retained. Its appearance makes explicit how collection changes the testing scale. The uniform amplitude clause supplies the distance guarantee at the amplitude named in \cref{obj:def:testing-handle}.

\subsection{From testing to squared estimation risk}

The estimation reduction uses generic priors over fixed schedules. Write \(\mathsf Q_+\) and \(\mathsf Q_-\), the complete-record mixture laws of \cref{obj:lem:full-record-two-prior-risk}, for the laws obtained by drawing a schedule from the respective prior and then independently running the canonical assignment--audit experiment. Thus, for a measurable record event, its mixture probability is the prior average of its conditional probability given the fixed schedule.

Suppose the positive and negative priors have constant causal targets \(+\Delta\) and \(-\Delta\), respectively, where \(\Delta>0\) is their common target magnitude. The reduction compares the overlap of the two record laws with the squared error required to estimate these separated targets. Its support condition connects prior-average loss to the minimax schedule class, and its measurability condition makes the induced record laws well defined. The following statement includes arbitrary measurable randomized estimators of the complete record.

\begin{lemmav}{lem:full-record-two-prior-risk}[Full-record two-prior risk bound]
Let \(n\) be the population size, \(d\) the degree bound, and \(q\) the audit-retention probability. Suppose:
\begin{itemize}
\item \textbf{(Parameters.)} \(n,d\) are integers satisfying \(n\ge4\) and \(1\le d\le n-1\), and \(q\in[0,1]\).
\item \textbf{(Priors.)} Two probability priors, each represented by a map from a measurable parameter space to fixed schedules, assign probability one to \(\mathcal M_{n,d}\) as defined in \cref{obj:def:model}.
\item \textbf{(Targets.)} For some \(\Delta>0\), the positive and negative priors respectively satisfy \(\tau(\theta)=\Delta\) and \(\tau(\theta)=-\Delta\) almost surely, where \(\tau(\theta)\), the population all-treated-versus-all-control contrast, is
\[
\tau(\theta)=\frac1n\sum_{i=1}^{n}
\bigl(Y_i(1,\ldots,1)-Y_i(0,\ldots,0)\bigr),
\]
and \(Y_i(z)\) denotes unit \(i\)'s additive potential outcome under schedule \(\theta\) at binary assignment \(z\).
\item \textbf{(Measurability.)} For each prior representation, the map from the parameter, assignment, and audit marks to the complete observed record is jointly measurable.
\end{itemize}
Under each prior, draw the schedule first, then independently draw \(Z\), the treatment-assignment vector, with independent Bernoulli-\(1/2\) coordinates, and \(W\), the array of audit marks on all ordered off-diagonal pairs, with independent Bernoulli-\(q\) coordinates, independently of \(Z\). The retained graph \(H\) and complete record \(O\) are
\[
H_{ij}=\mathbf1\{j\to i\text{ in }G\}\,W_{ij}\quad(i\ne j),
\qquad
O=(H,Z,Y(Z)),
\]
where \(G\) is the schedule's graph and \(Y(Z)\) is the vector of realized outcomes. Let \(\mathsf Q_+\) and \(\mathsf Q_-\) be the resulting complete-record mixture laws under the positive and negative priors. Then the minimax squared risk over measurable randomized estimators defined in \cref{obj:def:risk} satisfies
\[
R_{n,d}(q)\ge
\frac{\Delta^2}{2}
\left(1-\operatorname{TV}(\mathsf Q_+,\mathsf Q_-)\right),
\]
where \(\operatorname{TV}\) is the total-variation distance, defined as the supremum over measurable events of the absolute difference between their probabilities.
\end{lemmav}

The lower bound in \cref{obj:lem:full-record-two-prior-risk} combines target separation with record-law overlap. A larger constant target magnitude increases the squared-loss consequence of ambiguity, while a smaller total-variation distance preserves more overlap between the experiments. For the block priors, \cref{obj:lem:block-law} supplies support in the schedule class and the target magnitude \(mh/n\). The testing guarantees of \cref{obj:thm:block-testing-scale,obj:thm:nonlocal-testing-answer} therefore provide the record-law comparison needed by this reduction at the prescribed testing amplitude.

\subsection{Combining the comparisons across degrees and retention}

The block and supplied-graph comparisons address complementary parts of the admissible parameter range. The block experiment in \cref{obj:lem:block-law} requires room for equal numbers of sources and recipients. Its testing scale captures the dependence on association revelation, and its target magnitude records the fraction of the population occupied by recipients. The lower comparison in \cref{obj:lem:supplied-block-record-lower} applies throughout \(1\le d\le n-1\) and \(q\in[0,1]\), including degrees for which the population capacity restricts the block layout. It also identifies the degree-dependent estimation difficulty present when the true graph is supplied.

On the attainable side, \cref{obj:thm:observable-upper} controls the observable rule of \cref{obj:def:upper-rule}. The moment identity in \cref{obj:lem:score-audit} separates assignment variation from audit variation, and the upper envelope of \cref{obj:def:upper-envelope} governs the choice between the clipped score and zero. These ingredients supply the upper comparison for the named rule in \cref{obj:def:frontier-rule}; the block testing reduction and the supplied-graph lower comparison supply its minimax counterpart.

The risk scale \(F(n,d,q)\) of \cref{obj:def:frontier-scale} packages these comparisons through the source-association revelation probability. The positive-retention expression separates the degree and collection contributions and truncates their sum at one, representing saturation of the minimax risk order. For population sequences, these two contributions become the respective consistency criteria. The precision pair below records the already-defined scale and rule so that the explicit-constant statement can refer to them jointly.

\begin{definitionv}{def:frontier-handle}[Precision pair $(F,T^\star_{n,d,q})$]
The precision handle is the scale–rule pair \((F,T^\star_{n,d,q})\), defined by
\[
F(n,d,q)=
\begin{cases}
\min\left\{1,\dfrac{d^2}{n[1-(1-q)^d]}\right\},&q>0,\\
1,&q=0,
\end{cases}
\qquad
T^\star_{n,d,q}(O)=
\begin{cases}
\max\{-1,\min\{1,T_q(O)\}\},&q>0,\ u(n,d,q)<1,\\
0,&\text{otherwise}.
\end{cases}
\]
Here \(O\) is the complete observed record, \(T_q\) is the observable inverse-inclusion score, and \(u\) is the upper-risk envelope. The associated positive-retention elbow expression is
\[
\min\left\{1,\frac{d^2}{n}+\frac{d}{nq}\right\}.
\]
For admissible population sequences \(d_n,q_n\), the associated consistency criteria are
\[
\frac{d_n^2}{n}\longrightarrow0,
\qquad
\frac{d_n}{nq_n}\longrightarrow0,
\]
with the second ratio assigned \(+\infty\) when \(q_n=0\). The uniform risk comparison, equivalence with the elbow expression, and necessity and sufficiency of these criteria are proved properties of this pair. The original-record converse uses the complete conditional likelihood and retains the retained graph \(H\), assignment vector \(Z\), realized outcomes \(Y(Z)\), and observed global hidden treated-source count \(K\).
\end{definitionv}

The following statement supplies explicit universal constants for the primary frontier in \cref{obj:thm:observed-record-precision-frontier} and consolidates its endpoint and sequence deductions. The scale, rule, and estimator class remain those of \cref{obj:def:frontier-scale,obj:def:frontier-rule,obj:def:risk}.

\begin{theoremv}{thm:frontier-answer}[Explicit risk and consistency frontier]
Consider the canonical independent assignment--audit product design, with known edge-retention probability \(q\). Let \(\mathcal M_{n,d}\) be the fixed graph-and-coefficient schedules satisfying the common off-diagonal in-degree and out-degree bound \(d\) and the unit coefficient-mass bound. Write \(\mathcal L(T;\theta,q)\) for the expected squared error of \(T(O,\xi)\) in estimating \(\tau(\theta)\), the population all-treated-versus-all-control contrast, where \(O\) is the complete observed record and \(\xi\) is the independent uniform randomization seed. The minimax risk is
\[
R_{n,d}(q)
=
\inf_T\sup_{\theta\in\mathcal M_{n,d}}\mathcal L(T;\theta,q),
\]
where the infimum ranges over all measurable real-valued randomized estimators of the complete record.

Suppose that:
\begin{itemize}
\item \textbf{(Population size.)} \(n\) is a natural number with \(n\ge4\).
\item \textbf{(Degree bound.)} \(d\) is a natural number with \(1\le d\le n-1\).
\item \textbf{(Retention.)} \(q\in[0,1]\).
\end{itemize}
Then the risk scale \(F(n,d,q)\) and attaining rule \(T^\star_{n,d,q}\) of \cref{obj:def:frontier-scale,obj:def:frontier-rule} satisfy
\[
\frac{F(n,d,q)}{2560000\pi^2}
\le R_{n,d}(q)
\le
\sup_{\theta\in\mathcal M_{n,d}}
\mathcal L(T^\star_{n,d,q};\theta,q)
\le 24F(n,d,q).
\]

For every pair of sequences \(d_n\in\mathbb N\) and \(q_n\in\mathbb R\) satisfying
\[
1\le d_n\le n-1,\qquad q_n\in[0,1]
\qquad\text{for every }n\ge4,
\]
there exists a sequence of measurable randomized estimators \(T_n\) of the complete record such that
\[
\sup_{\theta\in\mathcal M_{n,d_n}}
\mathcal L(T_n;\theta,q_n)\longrightarrow0
\]
if and only if
\[
\frac{d_n^2}{n}\longrightarrow0
\qquad\text{and}\qquad
\begin{cases}
+\infty,&q_n=0,\\[2pt]
\dfrac{d_n}{nq_n},&q_n\ne0
\end{cases}
\longrightarrow0.
\]
All limits are as \(n\to\infty\).

For every finite population \(V\), every \(n,d\in\mathbb N\), every \(q\in\mathbb R\), and every record \(O\) on \(V\), the total observable rules of \cref{obj:def:frontier-rule,obj:def:upper-rule} agree:
\[
T^\star_{n,d,q}(O)=T^{\mathrm{up}}(O).
\]
The estimator attaining the displayed risk bound is the seeded measurable representative of this upper rule.

Finally, for every \(n\ge4\) and \(1\le d\le n-1\),
\[
F(n,d,0)=1,
\qquad
F(n,d,1)=\min\left\{1,\frac{d^2}{n}\right\}.
\]
For every \(q\in(0,1]\),
\[
\frac12\min\left\{1,\frac{d^2}{n}+\frac{d}{nq}\right\}
\le F(n,d,q)
\le
(1-e^{-1})^{-1}
\min\left\{1,\frac{d^2}{n}+\frac{d}{nq}\right\}.
\]
\end{theoremv}

The finite-population comparison in \cref{obj:thm:frontier-answer} gives matching risk scales with universal constants. The positive-retention expression makes the two contributions explicit: degree growth affects precision even with a supplied graph, and partial collection adds a retention-dependent contribution. At full retention the scale is \(\min\{1,d^2/n\}\); at zero retention it equals one. The attaining rule translates the upper comparison into an observable procedure using the known design parameters.

For admissible population sequences, sufficiency is constructive: the seeded measurable representatives of \(T^\star_{n,d_n,q_n}\) in \cref{obj:def:frontier-rule} attain uniformly vanishing squared risk whenever both displayed ratios tend to zero. Necessity applies to every sequence of measurable randomized estimators of the complete record: uniformly vanishing risk requires both the degree contribution and the collection contribution to vanish. Assigning the second ratio \(+\infty\) at zero retention incorporates that boundary into the same characterization.
